\documentclass[11pt]{article}
\usepackage{amsmath,amssymb,amsthm}
\newtheorem{proposition}{Proposition}
\title{An obstruction to a reflection-paired five-row local model}
\author{Research note; independently reviewed}
\date{30 September 2026}
\begin{document}
\maketitle
For a vector on five ordered rows write $\mathbb E$ for its coordinate
average, and let row reflection reverse the five coordinates. An even
vector is fixed by reflection; an odd vector changes sign.

\begin{proposition}
There is no rational local comparison model with centered columns
\[
 a=z+u,\quad b=z-u,\quad c=t+v,\quad d=t-v,
\]
where $z,t$ are odd and $u,v$ are even, satisfying the uniform interval
square-free moments through degree four:
\[
 \mathbb E\prod_{j\in S}y_j=
 \begin{cases}0&|S|\text{ odd},\\1/(|S|+1)&|S|\text{ even}.
 \end{cases}
\]
More generally, over the reals every such model has $u=0$ or $v=0$.
No monotonicity or interval constraint is needed for these assertions.
\end{proposition}
\begin{proof}
The first moments imply $\mathbb Eu=\mathbb Ev=0$. Since odd and
even vectors are orthogonal, the four mixed pair moments imply
$\mathbb Ezt=1/3$ and $\mathbb Euv=0$.
Suppose first that $u,v$ are nonzero. The space of even mean-zero
vectors on five rows has dimension two, so these orthogonal nonzero
vectors form a basis of that space.

Set $A=ab=z^2-u^2$ and $B=cd=t^2-v^2$. Both are even,
$\mathbb EA=\mathbb EB=1/3$, and the triple moments give
$\mathbb EAv=\mathbb EBu=0$ (their products with odd vectors have
zero mean automatically). Thus for some real $\alpha,\beta$,
\[
 A=\frac13+\alpha u,\qquad B=\frac13+\beta v.
\]
It follows that $\mathbb EAB=1/9$, contradicting the required fourth
moment $1/5$.

If $u=0$, then $a=b=z$ is an odd rational vector with
$\mathbb Ez^2=1/3$. Write $z=(r,s,0,-s,-r)$; this would require
$r^2+s^2=5/6$. The 3-adic valuation of a nonzero sum of two rational
squares is even: after removing a common power of three, the sum of
the two squared 3-adic integers is a unit. But $v_3(5/6)=-1$ is odd.
This is impossible. The case $v=0$ is identical.
\end{proof}
\end{document}
